\section{Experimental Realizations of the Framework}
\label{sec:experiments}

The preceding theory leaves several components deliberately modular.
The five studies below vary those components while keeping finite-codeword assignments as the common representation family; \cref{tab:experimental-framework-summary} summarizes the roles used in the experiments.
Code for the experimental studies is available at \url{https://github.com/yaniv-shulman/relational-compression}.

\begin{table}[t]
\centering
\scriptsize
\setlength{\tabcolsep}{2.5pt}
\renewcommand{\arraystretch}{1.08}
\setlength{\extrarowheight}{1pt}
\begin{tabularx}{\linewidth}{>{\raggedright\arraybackslash}p{0.13\linewidth}
                            >{\raggedright\arraybackslash}p{0.18\linewidth}
                            >{\raggedright\arraybackslash}p{0.25\linewidth}
                            >{\raggedright\arraybackslash}X}
\toprule
Role & Classical compression & Relational compression & Finite-codeword realization used here \\
\midrule
Source / fidelity-bearing content
& Source object $x$ and the content reconstructed from it
& Relational object $(V,S)$, with $S$ the fidelity-bearing structure conditional on shared $V$
& Explicit point geometry, graph structure, or teacher-defined relations; Study 4 instead uses source-image reconstruction to constrain the finite representation through a joint spatial decoder \\
Encoder-side observations
& Usually the source observation itself
& $X_V=(x_v)_{v\in V}$ when encoder-side observations are present
& Point coordinates, graph or node structural features, or image pixels mapped to finite assignments \\
Retained description
& Retained message or code $m$
& Retained relational description $m_S$
& The source-specific $m_S$ comprises declared finite assignments, codes, or tokens plus any required auxiliary payload; $q_\theta(\cdot\mid x)$ is an optimization parameterization where applicable \\
Reconstruction / evaluation
& Reconstructed source $\widehat x$
& $\widetilde S=\mathsf{Dec}_V(m_S)$, materialized or evaluated on demand
& Same-code equality or collision may supply the evaluated relation directly or feed a richer structural or reconstruction decoder \\
Primitive representation-side relation
& No universal analogue is prescribed
& Depends on the chosen description mechanism
& Hard equality $\ind[c_\theta(x)=c_\theta(x')]$ and soft collision $k_\theta(x,x')$ \\
Fidelity
& Distortion $d(x,\widehat x)$
& Relational distortion $\Delta_V(S,\widetilde S)$
& Study-specific source fidelity acting on or through collision, including centroid/pairwise, graph, normalized-cut, reconstruction-defined, and signed-teacher criteria \\
Resource / representation constraint
& Encoded bits, rate, alphabet size, or another declared resource
& Description cost $\Omega(m_S)$
& Alphabet bound $K$, binary width, or nominal latent capacity constrain the declared description convention but need not equal the complete resource $\Omega(m_S)$; any additional retained payload must also be counted \\
Realized organization
& Realized code distribution or entropy under the coding convention
& Separate from description resource unless explicitly incorporated into $\Omega$
& $H_2(q_Z)$, $K_{\mathrm{eff}}$, and $D_2(q_Z\|r)$ when used \\
Shared context
& Codebook or model information available to encoder and decoder
& Shared carrier $V$ and context $C$, including reusable model parameters
& Shared neural encoder or decoder parameters and fixed conventions not included in each source-specific retained description \\
\bottomrule
\end{tabularx}
\caption{Summary of the compression roles used to interpret the experimental realizations. The final column specializes the general relational-compression interface to the finite-codeword constructions used in this paper. Realized codeword organization is reported separately from operational coding rate unless the declared resource convention identifies them.}
\label{tab:experimental-framework-summary}
\end{table}

The five studies provide a compact progression through these roles.
Study 1 provides a controlled numerical realization of the exact centroid--relation correspondence and decoder decomposition; Study 2 holds the finite representation mechanism fixed while exchanging graph-relational fidelity criteria; and Study 3 carries the normalized-cut realization into a shared inductive encoder evaluated on unseen graphs.
Studies 4 and 5 then contrast reconstruction-defined requirements on a finite image representation with explicit teacher-defined relational requirements.\footnote{We write $\lambda_{\mathrm{org}}$ for marginal organization. In Studies 2--4 it corresponds to \texttt{lambda\_sep} (``separation weight'') in the implementation, with no change to settings or objectives; Study 5's pair-specific separation is distinct.}

\subsection{Study 1: A controlled reconstruction--relational identity}
\label{sec:synthetic-experiment}

The first study provides a controlled numerical realization of \cref{prop:centroid-pairwise}, checking that the implemented centroid and inverse-mass-weighted pairwise forms agree as scalar objectives and in their encoder gradients.
It also isolates what fails when the inverse-mass factor is omitted.

The source consists of 768 points in $\R^2$ from an imbalanced six-component Gaussian mixture.
A finite encoder has $K=8$ possible codewords, whose probabilities are evaluated exactly rather than sampled.
The mixture imbalance encourages nonuniform empirical codeword occupancy, making the inverse-mass normalization consequential.
Full settings and diagnostics are in \ref{app:synthetic-details}.

With $\hat q_Z(z):=\frac{1}{N}\sum_iq_{iz}$, the implemented pairwise objective is the finite-sample inverse-mass identity already given in \cref{eq:finite-soft-pairwise}:
\[
D_{\mathrm{pair}}
:=
\frac{1}{2N^2}
\sum_{z:\hat q_Z(z)>0}
\frac{1}{\hat q_Z(z)}
\sum_{i,j}
q_{iz}q_{jz}\|x_i-x_j\|_2^2.
\]
Thus $D_{\mathrm{pair}}=D_{\mathrm{cent}}$.
The deliberately uncorrected comparison is
\[
D_{\mathrm{raw}}
:=
\frac{1}{2N^2}
\sum_{i,j}
\|x_i-x_j\|_2^2
\sum_zq_{iz}q_{jz}.
\]
Uniform empirical occupancy would give $D_{\mathrm{cent}}=KD_{\mathrm{raw}}$.
The comparison therefore includes this global factor $K=8$; nonuniform occupancy instead requires the code-dependent inverse masses.

Across the optimization and diagnostic checks, the corrected pairwise and centroid objectives agree to approximately single-precision tolerance, with maximum discrepancies below $2\times10^{-6}$ and encoder-gradient cosines of 1.000.
By contrast, $KD_{\mathrm{raw}}$ has mean absolute error 4.79 in the random-state check.
A global scale does not replace the inverse aggregate-mass correction.

With the encoder assignments fixed, learned codeword reproduction vectors numerically realize \cref{eq:decoder-centroid-decomposition}.
The excess term $\sum_{z:\hat q_Z(z)>0}\hat q_Z(z)\|\bar x_z-\widehat x_z\|_2^2$ decreases from 21.05 to $8.15\times10^{-5}$, while the largest decomposition residual is $1.91\times10^{-6}$.
The result separates distortion due to the assignments from distortion due to an imperfect reproduction codebook.
The identity applies to per-codeword centroid reproduction; joint neural decoding of many latent tokens has a different structure.
\Cref{fig:synthetic-collision-centroid} summarizes the exact objective correspondence, the failure of the uncorrected collision control, and the decoder decomposition.

\begin{figure}[t]
\centering
\includegraphics[width=\linewidth]{synthetic_collision_centroid_equivalence.png}
\caption{Synthetic numerical realization of the centroid--pairwise correspondence. Top left: learned hard partition and soft centroids. Top right: the two exact objectives follow the same trajectory from identical initialization. Bottom left: corrected pairwise distortion lies on the centroid identity line, whereas $KD_{\mathrm{raw}}$ does not. Bottom right: decoder distortion decomposes into centroid distortion plus reproduction mismatch.}
\label{fig:synthetic-collision-centroid}
\end{figure}

The experiment verifies the inverse-mass centroid correspondence and separately exposes decoder reproduction mismatch, numerically illustrating why assignments and decoder reproductions are distinct parts of the description.

\subsection{Study 2: Exchanging relational fidelity on a fixed representation}
\label{sec:transductive-relational-distortion-experiment}

The second study demonstrates the framework's separation between the finite representation and the relational fidelity used to shape it.
Holding the graph-local $K=8$ representation and marginal-organization mechanism fixed, we exchange four source-derived graph fidelities and compare the resulting occupancy--fidelity behavior and partitions.
Each graph has a free $n\times K$ logit matrix with $K=8$, rather than a learned graph encoder.
There are no node features or labels, and no generalization problem.

For $q_i=\operatorname{softmax}(\ell_i/\tau)$ with $\tau=1$, all four source distortion criteria have the form
\begin{equation}
D_\rho(q)=\sum_{e=\{i,j\}\in E}\rho_e(1-q_i^\top q_j),\qquad
\rho_e\ge0,\quad\sum_{e\in E}\rho_e=1.
\label{eq:transductive-drho}
\end{equation}
The choices are normalized direct-edge distortion $D_E$ from \cref{eq:normalized-edge-distortion}, all-mode effective-resistance distortion $D_F$ from \cref{eq:source-normalized-dirichlet-distortion}, source-collision probe distortion $D_C$ from \cref{eq:source-collision-edge-distortion}, and entropy-field distortion $D_{H_2}$ from \cref{eq:entropy-probe-distortion} when its denominator is positive.
Here $D_{H_2}$ denotes distortion of the source transition-entropy field and is distinct from the code-occupancy statistic $H_2(\bar q_G)$ below.
The graph-local marginal is degree weighted, and each run minimizes
\begin{equation}
\mathcal L=D_\rho(q)+\lambda_{\mathrm{org}}D_2(\bar q_G\|U_K)
=D_\rho(q)-\lambda_{\mathrm{org}}H_2(\bar q_G)+\lambda_{\mathrm{org}}\log K.
\label{eq:transductive-relational-objective}
\end{equation}
Here $\bar q_G$ is the degree-weighted graph-local aggregate defined in \cref{eq:graph-aggregate-code}, and the second equality uses the uniform-reference R\'enyi identity in \cref{eq:uniform-renyi}.
Thus the sweep varies a noncollapse requirement while evaluating relational fidelity.

The source collections are 20 cleaned connected MalNet-Tiny test graphs, 20 cleaned COLLAB graphs, and 20 cleaned PROTEINS graphs~\citep{freitas2021large,morris2020tudataset,yanardag2015deep,borgwardt2005protein}.
To separate sensitivity to source edge weights from differences in topology, we also evaluate a weighted control on the same MalNet topologies, assigning deterministic positive lognormal edge weights.
Labels are ignored.
Because the graph-local objectives are optimized nonconvexly, every graph, fidelity, and organization weight uses the same multi-restart selection protocol, with selection by the target soft objective.
Full restart, refinement, graph-filtering, and sweep details are given in \ref{app:transductive-relational-distortion-details}.

\paragraph{Occupancy--fidelity trajectories}
At zero organization weight, every defined criterion selects the one-codeword solution, with zero positive-relation distortion and $K_{\mathrm{eff}}=1$.
Increasing organization moves the hard solutions toward broader codeword use with nonzero relational distortion.
On unweighted MalNet, $\lambda_{\mathrm{org}}=0.20$ yields hard $K_{\mathrm{eff}}\approx7.8$ across criteria, with own hard distortions 0.238 for $D_E$, 0.153 for $D_F$, 0.127 for $D_C$, and 0.096 for $D_{H_2}$.
These values quantify different source properties on their respective fidelity scales.

\begin{figure}[t]
\centering
\includegraphics[width=\linewidth]{transductive_relational_distortion_sweep.png}
\caption{Transductive collision-occupancy--fidelity trajectories. The horizontal coordinate is hard $H_2(\bar q_G)$ and the vertical coordinate is each criterion's own hard distortion. Points average the selected graphs at 15 organization weights. The secondary axis transforms the plotted mean entropy by exponentiation; it does not in general equal the arithmetic mean of per-graph effective counts.}
\label{fig:transductive-relational-distortion-sweep}
\end{figure}

In \cref{fig:transductive-relational-distortion-sweep}, complete collapse preserves all positive edges under the masking interpretation, while organization imposes distinctions that sever some of them, so the curves increase rather than follow the classical shape of an upper-rate distortion budget.
The experiment realizes the lower-occupancy-bound direction of \cref{eq:relational-complexity-distortion-frontier} through a penalty.
Its trajectories are empirical optimization results from a nonconvex solver rather than certified constrained frontiers.

\paragraph{Agreement between source criteria}
Some different fidelity constructions give similar edge priorities on particular sources.
\Cref{tab:transductive-distortion-correlations} compares their normalized importance vectors and their rankings of 256 approximately node-balanced random hard partitions per graph.
$D_F$ and $D_C$ are strongly aligned on unweighted MalNet and PROTEINS, with mean rank correlations 0.975 and 0.981.
The weighted MalNet control reduces the correlation to 0.724, showing that source weights materially affect their agreement.

\begin{table}[t]
\centering\small
\begin{tabular}{lrrrr}
\toprule
Source & \shortstack{Importance cosine\\$(D_F,D_C)$} & \shortstack{Spearman\\$(D_F,D_C)$} & \shortstack{Spearman\\$(D_F,D_{H_2})$} & \shortstack{Spearman\\$(D_C,D_{H_2})$} \\
\midrule
MalNet unweighted & 0.979 & 0.975 & 0.666 & 0.616 \\
MalNet weighted & 0.744 & 0.724 & 0.500 & 0.550 \\
COLLAB & 0.833 & 0.814 & 0.523 & 0.774 \\
PROTEINS & 0.983 & 0.981 & 0.616 & 0.608 \\
\bottomrule
\end{tabular}
\caption{Source-fidelity agreement. Importance cosine compares the normalized source-importance vectors for $D_F$ and $D_C$; the remaining columns report Spearman correlations of distortion values over random hard partitions, averaged over graphs. COLLAB correlations involving $D_{H_2}$ use the 18 sources where that normalized fidelity is defined.}
\label{tab:transductive-distortion-correlations}
\end{table}

The entropy-field criterion is generally less aligned with the other probes.
On PROTEINS at $\lambda_{\mathrm{org}}=0.20$, its solutions have hard $K_{\mathrm{eff}}\approx5.59$ and own distortion 0.091, while the other criteria have effective counts near 4.2 and own distortions 0.149--0.163.
Because these differences mix occupancy and fidelity, the occupancy-matched analysis below is needed to compare preserved structure at similar complexity.
COLLAB also has a slower transition, reaching mean hard $K_{\mathrm{eff}}\ge4$ only at weight 0.50, with less sharp soft assignments at high weights.

\paragraph{Occupancy-matched partitions and transfer}
To separate those effects, solutions are selected independently from the existing sweeps at hard $H_2$ targets from 0.25 to 2.0 nats.
A pair is compared only when each solution is within 0.10 nats of the target and their entropies differ by at most 0.10 nats.
Adjusted Rand index (ARI) evaluates label-invariant partition agreement.
Cross-distortion excess uses one evaluation fidelity at a time: $\Delta_{A\to B}=D_B(c_A)-D_B(c_B)$ and $\Delta_{B\to A}=D_A(c_B)-D_A(c_A)$.
These are approximate occupancy matches of heuristic solutions; the reported excess is descriptive and can be slightly negative.
\Cref{fig:transductive-rate-matched-partition-agreement} summarizes the resulting hard-partition agreement across occupancy targets and criterion pairs.

\begin{figure}[t]
\centering
\includegraphics[width=\linewidth]{transductive_rate_matched_partition_agreement.png}
\caption{Occupancy-matched hard-partition agreement. Mean ARI is shown with interquartile bands for criterion pairs meeting the 0.10-nat matching rule; a point is displayed only when at least five graphs match. Missing low-entropy MalNet points reflect limited sweep coverage.}
\label{fig:transductive-rate-matched-partition-agreement}
\end{figure}

On PROTEINS, $D_F$ and $D_C$ produce mean ARI 0.968 at $H_2=1.0$ and 0.991 at $H_2=1.25$, with near-zero transfer excess.
At $H_2=2.0$, their ARI is 0.764 while the two excesses remain only 0.0046 and 0.0084.
Thus partitions can differ without substantially changing either criterion's preserved content.
By contrast, at $H_2=2.0$ on unweighted MalNet, $D_F$ versus $D_{H_2}$ has ARI 0.166 and $D_C$ versus $D_{H_2}$ has ARI 0.196, compared with 0.781 for $D_F$ versus $D_C$.
The full matched table, six-pair diagnostics, and transfer matrix are retained in \ref{app:transductive-relational-distortion-details}.

Thus, with representation and marginal organization held fixed, changing source fidelity changes source-edge priorities and can change the resulting partition; agreement in priorities, assignments, and cross-fidelity transfer need not coincide.

\subsection{Study 3: An inductive realization of normalized-cut fidelity}
\label{sec:malnet-ncut-experiment}

The third study moves from graph-specific optimization to a shared encoder evaluated on unseen graphs.
It demonstrates an inductive realization of the normalized-cut correspondence: a single finite-codeword encoder is trained with graph-local soft normalized cut while marginal organization is varied, and the resulting hard partitions are evaluated by normalized-cut fidelity, code occupancy, and the post-hoc fidelity $D_F$.
The numerical comparison below uses a per-graph spectral normalized-cut reference.
For related differentiable graph-partitioning constructions, we refer the reader to GAP, the method of Gatti et al., and MinCutPool~\citep{nazi2019gap,gatti2022deep,bianchi2020spectral}.
The probabilistic-cut work of \citet{ghriss2026beyond} addresses expected discrete cut objectives with random denominators, whereas this study uses the deterministic soft-mass surrogate.

\paragraph{Source objects and preprocessing}
MalNet-Tiny function-call graphs are split using the official train, validation, and test assignments~\citep{freitas2021large}.
The inputs, losses, and selection criteria exclude malware labels.
Each directed graph is converted to a simple undirected edge union, self-loops and duplicate edges are removed, and the largest connected component is reindexed.
Requiring at least 512 nodes after cleaning leaves 1995 training, 275 validation, and 584 test graphs.
The exclusions are documented in \ref{app:malnet-ncut-details}.

Node features are structural.
On the cleaned graph they include a constant, degree, $\log(1+d_i)$, graph-normalized degree, and random-walk return estimates for $\operatorname{diag}(P^t)$, where $P$ is the transition matrix of the cleaned undirected graph.
The estimates are cached deterministically for each graph.
Exact return probabilities are permutation-equivariant node features; the finite cached probe realization retains exact samplewise equivariance when its probes are transported with a relabeling.

Directed in-degree, out-degree, total degree, their logarithmic and normalized forms, and PageRank relative to the uniform baseline are computed on the original directed graph before symmetrization and carried to the retained vertices.
Explicit node identifiers and Laplacian eigenvectors are also excluded from the model inputs.
The encoder can use directed structural context even though fidelity is evaluated on the cleaned undirected graph.
This available-input convention is part of the inductive realization.

\paragraph{Encoder and objective}
A GraphGPS-style encoder~\citep{rampasek2022recipe} combines local GIN message passing~\citep{xu2019powerful} with global factorized efficient attention~\citep{shen2021efficient}, followed by a graph-context-conditioned $K=8$ assignment head.
Full architecture and training settings are given in \ref{app:malnet-ncut-details}.
Soft assignments are $q_i=\operatorname{softmax}(\ell_i)$; hard assignments take the maximizing codeword.

For each graph, $\bar q_G$ is the degree-weighted graph-local aggregate defined in \cref{eq:graph-aggregate-code}.
The minibatch loss is
\begin{equation}
\mathcal L_B=\frac1{|B|}\sum_{G\in B}
\left[\operatorname{Ncut}_{K,\mathrm{soft}}(G)
+\lambda_{\mathrm{org}}D_2(\bar q_G\|U_K)\right].
\label{eq:malnet-ncut-loss}
\end{equation}
Here $\operatorname{Ncut}_{K,\mathrm{soft}}$ is the fixed-$K$ soft normalized-cut criterion of \cref{eq:soft-ncut-relational-distortion}, while $D_2(\bar q_G\|U_K)$ uses the uniform-reference organization term of \cref{eq:uniform-renyi}.
The graph-local aggregate supplies the endogenous inverse-mass construction of \cref{eq:soft-nassoc-affinity}, while the external uniform reference specifies marginal organization.
The loss is averaged over graphs only after these graph-local quantities are formed.
When $\lambda_{\mathrm{org}}=0$, the marginal organization term is absent, so training minimizes only the graph-local soft normalized-cut criterion.
Training uses five seeds.

\paragraph{Evaluation}
Validation mean hard fixed-$K$ normalized cut, following \cref{eq:soft-ncut-relational-distortion,eq:hard-ncut}, selects checkpoints.
An empty hard class contributes zero association and hence a unit penalty, so using fewer than eight classes receives the corresponding fixed-$K$ cost.
Hard evaluation uses the induced partition relation against the cleaned source graph.
Post-hoc $D_F$ from \cref{eq:hard-partition-dirichlet-distortion} evaluates a second relational property after training.

A spectral reference computes a normalized-Laplacian embedding and deterministic $K$-means discretization separately for each cleaned graph, following the standard spectral-clustering construction~\citep{ng2001spectral}, and provides a transductive objective comparison.
The learned model instead applies one parameter set after structural preprocessing of each held-out graph.

\begin{table}[t]
\centering\small
\begin{tabular}{lrrrr}
\toprule
$\lambda_{\mathrm{org}}$ & Val Ncut & Test Ncut & Test $D_F$ & Test median Ncut \\
\midrule
0 & $1.129\pm0.063$ & $1.160\pm0.043$ & $0.0606\pm0.0115$ & $1.083\pm0.043$ \\
0.03 & $1.137\pm0.060$ & $1.171\pm0.062$ & $0.0657\pm0.0179$ & $1.093\pm0.061$ \\
0.10 & $1.158\pm0.029$ & $1.186\pm0.033$ & $0.0926\pm0.0055$ & $1.116\pm0.031$ \\
0.20 & $1.136\pm0.039$ & $1.170\pm0.030$ & $0.1117\pm0.0017$ & $1.114\pm0.035$ \\
0.50 & $1.285\pm0.026$ & $1.319\pm0.022$ & $0.1475\pm0.0020$ & $1.270\pm0.017$ \\
\midrule
Spectral reference & 1.143 & 1.159 & 0.0805 & 1.104 \\
\bottomrule
\end{tabular}
\caption{MalNet-Tiny fidelity results. Learned rows show mean and sample standard deviation across five seeds. Lower is better for all displayed fidelities. The deterministic spectral reference is evaluated once on the same cleaned splits. $D_F$ is post-hoc.}
\label{tab:malnet-ncut-results}
\end{table}

The organization diagnostics report soft $K_{\mathrm{eff}}$, the collision-effective code count defined in \cref{eq:effective-code-count}, and soft $D_2$, meaning $D_2(\bar q_G\|U_K)$ with the uniform-reference identity in \cref{eq:uniform-renyi}.

\begin{table}[t]
\centering\small
\begin{tabular}{lrrr}
\toprule
$\lambda_{\mathrm{org}}$ & Largest hard volume fraction & Soft $K_{\mathrm{eff}}$ & Soft $D_2$ \\
\midrule
0 & $0.840\pm0.042$ & $1.42\pm0.14$ & $1.745\pm0.093$ \\
0.03 & $0.807\pm0.072$ & $1.54\pm0.26$ & $1.679\pm0.149$ \\
0.10 & $0.678\pm0.023$ & $2.00\pm0.08$ & $1.416\pm0.043$ \\
0.20 & $0.544\pm0.006$ & $2.45\pm0.03$ & $1.195\pm0.012$ \\
0.50 & $0.375\pm0.012$ & $3.55\pm0.10$ & $0.820\pm0.030$ \\
\midrule
Spectral reference & 0.575 & -- & -- \\
\bottomrule
\end{tabular}
\caption{Organization of the same MalNet-Tiny solutions. Quantities are computed per graph before averaging. Mean effective count and mean divergence need not obey a nonlinear identity after averaging. The reference has a hard effective count of 2.78; it has no soft-assignment counterpart in these columns.}
\label{tab:malnet-organization-results}
\end{table}

\paragraph{Fidelity and organization}
Pure normalized cut attains test hard Ncut $1.160\pm0.043$, close to the spectral mean 1.159.
The learned partition is more concentrated: its largest hard class contains mean volume fraction $0.840\pm0.042$, and soft $K_{\mathrm{eff}}$ is only $1.42\pm0.14$ despite nearly all eight hard classes being active.
The two solutions therefore have similar objective values but different representation organization.

Increasing organization weight to 0.20 reduces the largest hard volume fraction to $0.544\pm0.006$ and raises soft $K_{\mathrm{eff}}$ to $2.45\pm0.03$, while test Ncut remains $1.170\pm0.030$.
At weight 0.50, the volume fraction falls to $0.375\pm0.012$ and effective count rises to $3.55\pm0.10$, but Ncut increases to $1.319\pm0.022$.
The occupancy trend is stable across the sweep, while small Ncut differences among moderate weights should not be overinterpreted relative to seed variation.
\Cref{fig:malnet-ncut-sweep} summarizes this occupancy--fidelity behavior across the organization sweep.

\begin{figure}[t]
\centering
\includegraphics[width=\linewidth]{malnet_ncut_separation_sweep.png}
\caption{MalNet-Tiny organization sweep. Error bars show sample standard deviations over five seeds. Increasing the marginal-organization weight reduces dominant-class volume and broadens soft effective occupancy; post-hoc $D_F$ shows the accompanying change in another relational fidelity. Dashed lines show transductive spectral summaries.}
\label{fig:malnet-ncut-sweep}
\end{figure}

The post-hoc probe also distinguishes partitions that look similar under normalized cut.
Weights 0.03 and 0.20 give almost identical test Ncut, $1.171\pm0.062$ and $1.170\pm0.030$, but $D_F$ values $0.0657\pm0.0179$ and $0.1117\pm0.0017$.
After averaging seeds and pairing by test graph, the latter has higher $D_F$ on 581 of 584 graphs, with mean difference 0.0460; the corresponding mean Ncut difference is $-0.0006$.
Normalized cut and effective-resistance distortion therefore define different fidelities.

Likewise, the pure learned solution has lower mean $D_F$ than the spectral reference, 0.0606 versus 0.0805, but is substantially more concentrated.
The comparison matches nominal assignment capacity but not realized occupancy; the learned and spectral procedures also have different deployment and model-cost conventions.
\Cref{fig:malnet-ncut-partition-example} illustrates cases near the 25th, 50th, 75th, and 95th percentiles of the learned-minus-spectral normalized-cut gap among held-out graphs with at most 900 nodes.

\begin{figure}[t]
\centering
\includegraphics[width=\linewidth]{malnet_ncut_partition_example.png}
\caption{Representative MalNet-Tiny partitions at organization weight 0.20, selected near the 25th, 50th, 75th, and 95th percentiles of the learned-minus-spectral hard Ncut gap among held-out graphs with at most 900 nodes. Each row shares a force-directed layout. Partition labels and color assignments are independently permuted for display, so colors do not correspond across panels.}
\label{fig:malnet-ncut-partition-example}
\end{figure}

Together, the study shows that a shared encoder can learn the normalized-cut realization on unseen graphs, while occupancy and post-hoc $D_F$ reveal distinctions not captured by normalized cut alone.

\subsection{Study 4: Reconstruction-defined image representations}
\label{sec:flowers102-experiment}

The fourth study demonstrates a different route by which a source requirement can shape a finite-codeword representation: pixel reconstruction acts through a joint spatial decoder rather than directly on same-code relations.
Holding the 16-bit bottleneck architecture fixed, we vary marginal organization to examine how reconstruction quality and hard-code occupancy interact; a 32-bit run provides a separate capacity check.
Unlike the analytical centroid realization, the synthesis transform jointly interprets neighboring spatial tokens, so no exact per-codeword centroid identity is assumed.
Reconstruction MSE is the source fidelity, while hard collision entropy and effective code count describe representation organization.

Binary and discrete image autoencoders have substantial precedents.
Recurrent binary bottlenecks, VQ-VAE, soft-to-hard vector quantization, and learned transform codecs already combine reconstruction with constrained latent representations~\citep{toderici2016variable,vanDenOord2017neural,agustsson2017soft,balle2017end,balle2018variational}.
Related variational representation-learning work explicitly studies the rate--distortion trade-off between compression and reconstruction accuracy~\citep{alemi2018fixing}.
The Latent Bernoulli Autoencoder is a particularly close precedent: it trains an end-to-end reconstruction autoencoder with a hard $\{-1,+1\}$ latent obtained by zero thresholding and a straight-through gradient surrogate~\citep{fajtl2020latent}.
Codebook-free alternatives include finite scalar quantization, lookup-free quantization, and binary spherical quantization~\citep{mentzer2024finite,yu2024language,zhao2025binary}, while token-prior regularization also precedes this study~\citep{zhang2023regularized}.
The experiment isolates a collision-based organization mechanism in a purpose-limited architecture.

\paragraph{Data and representation}
Oxford Flowers102 contains 8189 images from 102 categories~\citep{nilsback2008automated}; labels are not used.
For this reconstruction study, the official test split supplies the 6149 training images, the official validation split supplies 1020 validation images, and the official training split supplies 1020 test images.
This nonstandard split is reported explicitly as part of the reconstruction protocol.
Training uses image augmentation; validation and test use deterministic $256\times256$ center-crop preprocessing.

The convolutional analysis transform produces $b$ logits at each position of a $32\times32$ latent grid.
Each position is a spatial token with $b$ hard bits, and code statistics pool tokens across images and positions.
The nominal raw latent capacity is $b/64$ bits per input pixel: 0.25 bpp for 16 bits and 0.50 bpp for 32 bits.
These are fixed latent capacities before any entropy coding.

The convolutional synthesis transform jointly interprets the spatial tokens to produce a sigmoid RGB output.
GDN/IGDN follows established nonlinear transform-coding practice~\citep{balle2017end}.
There are no encoder--decoder skip connections, so every reconstruction passes through the bottleneck.
The encoder and decoder are shared across images; the per-image nominal latent bpp counts the latent bits, while model size belongs to the complete coding convention.

At evaluation, bit $r$ at position $u$ is $h_r(x,u)=1$ if its logit $\ell_r(x,u)>0$, and $-1$ otherwise.
Hard tokens therefore use the centered sign convention $\{-1,+1\}^b$, a relabeling of the binary alphabet $\{0,1\}^b$ in \cref{eq:factorized-code}.
Centering the two states symmetrically about the zero threshold is convenient for the relaxation through which gradients are propagated during training, while hard equality and collision statistics are unchanged by the relabeling.
Training uses the mean-group relaxation described in \ref{app:flowers102-details}.
DiffPrune provides earlier work on deterministic approximate binary gates~\citep{shulman2020diffprune}.
The training surrogate used here more directly adapts the sign-partitioned group-mean transformation of \citet{shulman2021exact}, originally introduced for binary network parameters, to spatial latent variables, and combines it with the concentration penalty of \cref{eq:concentration-reduction} rather than reproducing the original continuation formulation unchanged.
Binary latent autoencoders and sign-based discrete tokenizers precede this study; to our knowledge, prior discrete-autoencoder work has not used this group-coupled relaxation to train a hard-sign latent bottleneck.

\paragraph{Collision-based organization}
The training loss has the form
\begin{equation}
\mathcal L=\mathcal L_{\mathrm{mse}}
+\lambda_{\mathrm{conc}}\mathcal L_{\mathrm{conc}}
+\lambda_{\mathrm{org}}\widetilde D_2.
\label{eq:flowers-objective}
\end{equation}
The last term acts on a differentiable threshold-margin surrogate, while evaluation uses literal collisions of the hard sign tokens.
Let $s_r>0$ be a detached per-channel RMS logit scale maintained by an exponential moving average.
For gain $\gamma>0$, define
\begin{equation}
\eta_r(x,u):=\ell_r(x,u)/s_r,\qquad
\tilde p_r(x,u)=\sigma(\gamma\eta_r(x,u)).
\label{eq:flowers-threshold-probability}
\end{equation}
Equivalently, an independent logistic threshold perturbation $T_r$ with location zero and scale $s_r/\gamma$ gives $\tilde p_r=\Pr[T_r<\ell_r]$.
Independence is across bit draws and independently encoded tokens; using one shared random threshold would define a different collision model.

For two tokens, whole-code collision under these factorized probabilities is
\begin{equation}
\widetilde k_{ij}
=\prod_{r=1}^b\left[\tilde p_{ir}\tilde p_{jr}+(1-\tilde p_{ir})(1-\tilde p_{jr})\right].
\label{eq:flowers-threshold-collision}
\end{equation}
This is the factorized binary collision of \cref{eq:bit-collision}, instantiated for the code in \cref{eq:factorized-code} with the smooth threshold probabilities $\tilde p_r$.
The sampled organization loss is
\begin{equation}
\widetilde D_2=b\log2+\log\widehat{\E}[\widetilde k_{ij}],
\label{eq:flowers-separation}
\end{equation}
Since $K=2^b$, the population identities in \cref{eq:population-collision,eq:uniform-renyi} give $D_2(q_Z\|U_K)=b\log2+\log\E[k_\theta(X,X')]$.
The implemented $\widetilde D_2$ replaces that population collision by the smooth threshold model and a finite minibatch pair estimate, so it is an optimization surrogate rather than the hard-code $H_2$ and $K_{\mathrm{eff}}$ statistics reported at evaluation.
Because the logarithm is applied to a finite sampled collision estimate, $\widetilde D_2$ is not generally an unbiased estimator of the corresponding population R\'enyi divergence, nor is a finite-sample realization guaranteed to inherit its nonnegativity.
The reported hard $H_2$ and $K_{\mathrm{eff}}$ are computed from the empirical distribution of complete hard spatial tokens pooled across images and positions in the evaluation split, using \cref{eq:collision-entropy,eq:effective-code-count}; this pooled token collection is the declared evaluation population.

\paragraph{Ablations and results}
The 16-bit sweep uses weights $0$, $10^{-7}$, $10^{-6}$, and $10^{-5}$, plus a 32-bit capacity run at $10^{-6}$.
All settings use the same data split and one seed, 1337; validation hard MSE selects checkpoints.
Large changes therefore illustrate mechanism while small differences remain statistically unresolved.
Full checkpoint, MSE, MS-SSIM, and optimization settings are retained in \ref{app:flowers102-details}.

\begin{table}[t]
\centering\small
\begin{tabular}{lrrrrr}
\toprule
Setting & Bits & Nominal bpp & PSNR & Hard $H_2$ & $K_{\mathrm{eff}}$ \\
\midrule
No organization & 16 & 0.25 & 19.86 & 1.98 & 7.2 \\
Weaker, $10^{-7}$ & 16 & 0.25 & 24.36 & 7.00 & 1\,098 \\
Baseline, $10^{-6}$ & 16 & 0.25 & 24.83 & 9.50 & 13\,388 \\
Stronger, $10^{-5}$ & 16 & 0.25 & 24.56 & 10.29 & 29\,569 \\
Capacity, $10^{-6}$ & 32 & 0.50 & 25.88 & 12.46 & 257\,753 \\
\bottomrule
\end{tabular}
\caption{Flowers102 hard-code reconstruction and organization at validation-selected checkpoints. PSNR is in dB and $H_2$ in nats. Each row is one run. The nominal bpp describes raw latent bits before entropy coding.}
\label{tab:flowers102-results}
\end{table}

Without organization, reconstruction learns a nontrivial partition but uses only 38 distinct codewords, with effective count about 7.2, hard collision $1.38\times10^{-1}$, and 19.86 dB PSNR.
Thus reconstruction supplies a reason to distinguish inputs while allowing concentrated use of the available 16-bit alphabet.
Several bit channels become saturated in this run.

A weight of $10^{-7}$ raises PSNR to 24.36 dB and effective count to about $1.1\times10^3$.
The $10^{-6}$ setting gives the best 16-bit reconstruction among these runs, 24.83 dB and 0.896 MS-SSIM, with effective count 13\,388.
Increasing the weight to $10^{-5}$ further broadens occupancy to 29\,569 effective codes while PSNR decreases to 24.56 dB.
The few-tenths-of-a-decibel difference is unresolved in this single-seed study, while the sweep suggests that utilization can continue increasing after reconstruction quality plateaus.
\Cref{fig:flowers102-dynamics} shows how reconstruction and hard-code utilization evolve across the organization sweep.
\begin{figure}[t]
\centering
\includegraphics[width=\linewidth]{flowers102_ablation_dynamics.png}
\caption{Validation dynamics of the image reconstruction ablations. Utilization is measured from hard spatial tokens; smooth probabilities supply the training surrogate. The zero-weight run is highly concentrated; larger weights broaden occupancy. The 32-bit run has a larger nominal bottleneck capacity.}
\label{fig:flowers102-dynamics}
\end{figure}

Support size and effective occupancy tell different stories.
The baseline and stronger 16-bit settings activate 65\,208 and 65\,528 of 65\,536 possible words, but their effective counts are only 13\,388 and 29\,569.
The 32-bit run improves reconstruction to 25.88 dB and 0.920 MS-SSIM with effective count about $2.58\times10^5$.
It serves as a capacity check in which the nominal bpp, code space, and uniform-reference target all change.
Empirical occupancy also remains sample dependent, especially in a very large alphabet.
\Cref{fig:flowers102-16bit-tradeoff} summarizes the four fixed-capacity 16-bit operating points.
\begin{figure}[t]
\centering
\includegraphics[width=0.72\linewidth]{flowers102_16bit_tradeoff.png}
\caption{The four 16-bit test operating points. Organization initially improves both effective code use and PSNR; beyond the baseline it increases utilization while reconstruction slightly worsens. The single-seed operating points provide a descriptive trajectory rather than a statistically resolved frontier.}
\label{fig:flowers102-16bit-tradeoff}
\end{figure}

Qualitative hard-code reconstructions are retained in \cref{fig:flowers102-reconstruction-examples}.
Taken together, reconstruction can train a hard finite representation through a joint decoder without ensuring broad alphabet use; marginal organization changes that use separately, while token collision remains an organization statistic rather than the reconstruction fidelity itself.
The 32-bit run separately probes increased nominal capacity.

\subsection{Study 5: Teacher-defined relations without reconstruction}
\label{sec:flowers102-teacher-experiment}

The fifth study asks whether an externally supplied source relation can organize a finite hard representation without reconstruction or a marginal-organization penalty.
A frozen visual teacher defines baseline-relative favored and disfavored relations among spatial patches, and an independent student encoder is trained only through pair-specific whole-code collision.
Hard relation-conditioned collision provides the primary evaluation.
A decoder fitted only after freezing the encoder supplies a secondary cross-fidelity probe of recoverable image structure.

Transferring relations rather than individual features is established in relational and similarity-preserving distillation~\citep{park2019relational,tung2019similarity}.
Dense visual precedents include STEGO's distillation of feature correspondences, TokenCut's use of transformer patch affinities, and DeepCut's graph clustering of pretrained image features~\citep{hamilton2022unsupervised,wang2022selfsupervised,aflalo2023deepcut}.
Vector-quantized teacher features also appear in BEiT v2~\citep{peng2022beitv2}.
The selected realization combines an independent image encoder, cross-image baseline-relative relations, and a whole-code collision objective.

\paragraph{Source relation and student}
The data split and image preprocessing match those of \cref{sec:flowers102-experiment}.
A frozen DINO ViT-S/8 teacher~\citep{caron2021emerging} sees the same $256\times256$ image crop as the student, with the teacher's input normalization.
Its final normalized spatial tokens form a $32\times32$ grid, matching the student's latent grid.
Let $e_i$ be the frozen teacher embedding of spatial patch token $i$.

The student uses the same convolutional analysis transform as Study 4, with $b=16$ logits per spatial position.
Hard evaluation uses zero-threshold signs.
During representation learning,
\begin{equation}
p_r(i)=\sigma(\ell_r(i)/\tau_c),\qquad \tau_c=0.25.
\end{equation}
There is no mean-group relaxation, concentration term, marginal organization penalty, entropy model, or reconstruction decoder at this stage.

\paragraph{One teacher relation, two relational roles}
Each minibatch supplies a bounded sample of spatial tokens, and valid candidates for token $i$ form $\mathcal N_i$.
Self-pairs and same-image pairs are excluded in the reported run.
For teacher temperature $\tau_T=0.1$, define
\begin{equation}
S^{\mathrm{patch}}_{ij}
=\frac{\exp(e_i^\top e_j/\tau_T)}
{\sum_{j'\in\mathcal N_i}\exp(e_i^\top e_{j'}/\tau_T)},
\qquad
S_{\mathrm{bg},ij}=\nu_i=|\mathcal N_i|^{-1}.
\label{eq:teacher-row-relation}
\end{equation}
The teacher mapping and relation-construction rule are fixed, but $S^{\mathrm{patch}}$ is row-normalized within each sampled candidate set and is therefore not one immutable affinity matrix over the full dataset.
The source-side baseline is uniform over valid candidates.
The resulting favored and disfavored pair weights are detached constants during optimization:
\begin{equation}
S_{+,ij}=[S^{\mathrm{patch}}_{ij}-S_{\mathrm{bg},ij}]_+,\qquad
S_{-,ij}=[S_{\mathrm{bg},ij}-S^{\mathrm{patch}}_{ij}]_+.
\label{eq:teacher-signed-weights}
\end{equation}
Thus a single teacher relation supplies both favored and disfavored collisions through the baseline-relative construction of \cref{eq:signed-relation}.

The cosine and candidate mask are symmetric, but row normalization makes $S^{\mathrm{patch}}$ directional.
The student collision is symmetric and both directions are included.
Accordingly, each side of the loss aggregates the corresponding directional weight masses on an unordered pair.
This aggregation is performed separately for favored and disfavored weights and can differ from thresholding a single symmetrized signed relation.

The differentiable whole-code probability is
\begin{equation}
\widetilde k_{ij}=\prod_{r=1}^b[p_r(i)p_r(j)+(1-p_r(i))(1-p_r(j))].
\end{equation}
This is the factorized whole-code collision of \cref{eq:bit-collision}, evaluated using the student's smooth Bernoulli probabilities.
The representation-side neutral collision level is the population collision induced by the uniform codeword reference $U_K$:
\[
\sum_z U_K(z)^2=\frac{1}{K}=2^{-b}.
\]
Define
\begin{equation}
\delta_{ij}:=\operatorname{logit}(\widetilde k_{ij})-\operatorname{logit}(1/K).
\label{eq:teacher-centered-margin}
\end{equation}
The implementation clips the probability away from zero and one before evaluating the logit; the numerical convention is described in \ref{app:flowers102-teacher-details}.
The source baseline $S_{\mathrm{bg},ij}$ determines the favored and disfavored requirements, while $1/K$ centers the representation-side margin.
The uniform distribution $U_K$ is used here only to define this neutral whole-code collision level: the objective does not match the aggregate distribution $q_Z$ to $U_K$ and has no $\mathcal L_{\mathrm{org}}$ term.

The two sides of the objective are
\begin{equation}
\begin{aligned}
\mathcal L_{\mathrm{align}}
&=\frac{\sum_{ij}S_{+,ij}\operatorname{softplus}(-\delta_{ij})}{\sum_{ij}S_{+,ij}},\\
\mathcal L_{\mathrm{sep}}^{\mathrm{rel}}
&=\frac{\sum_{ij}S_{-,ij}\operatorname{softplus}(\delta_{ij})}{\sum_{ij}S_{-,ij}},\\
\mathcal L_{\mathrm{teacher}}
&=\tfrac12\left(\mathcal L_{\mathrm{align}}+\mathcal L_{\mathrm{sep}}^{\mathrm{rel}}\right),
\end{aligned}
\label{eq:dino-teacher-loss}
\end{equation}
with zero-mass sides omitted from the average and zero loss when neither side is active.
Favored relations encourage collision above the reference, while disfavored relations encourage it below the reference.
The objective preserves baseline-relative distinctions without reconstructing teacher cosine values.

\paragraph{Evaluation and results}
The representation checkpoint is selected by validation teacher loss.
Hard sign tokens provide the reported occupancy statistics and relation-conditioned collision rates.
Let $c(i)\in\mathcal Z$ denote the complete hard codeword assigned to patch token $i$ by zero-thresholding the student logits.
The relation-conditioned collision rates are
\[
\widehat\kappa_+
=\frac{\sum_{ij}S_{+,ij}\ind[c(i)=c(j)]}{\sum_{ij}S_{+,ij}},
\qquad
\widehat\kappa_-
=\frac{\sum_{ij}S_{-,ij}\ind[c(i)=c(j)]}{\sum_{ij}S_{-,ij}}.
\]
The pooled empirical hard-code collision is
\[
\widehat\kappa_{\mathrm{pop}}
=\sum_z\widehat q_Z(z)^2
=e^{-\widehat H_2}.
\]
Favored-pair enrichment and disfavored-pair suppression are, respectively,
\[
\frac{\widehat\kappa_+}{\widehat\kappa_{\mathrm{pop}}},
\qquad
1-\frac{\widehat\kappa_-}{\widehat\kappa_{\mathrm{pop}}}.
\]
These ratios use the pooled empirical hard-code collision as their baseline and are descriptive rather than calibrated against a candidate-matched null.

\begin{table}[t]
\centering\small
\begin{tabular}{lrr}
\toprule
Metric & Validation & Test \\
\midrule
Representation epoch & \multicolumn{2}{c}{146} \\
Teacher loss & 0.389 & 0.389 \\
Favored-pair hard collision & $1.35\times10^{-2}$ & $1.29\times10^{-2}$ \\
Disfavored-pair hard collision & $1.37\times10^{-4}$ & $1.47\times10^{-4}$ \\
Aggregate hard collision & $1.03\times10^{-3}$ & $1.04\times10^{-3}$ \\
Favored-pair enrichment & 13.13 & 12.40 \\
Disfavored-pair suppression & 0.866 & 0.858 \\
Hard $H_2$ (nats) & 6.88 & 6.87 \\
$K_{\mathrm{eff}}$ & 974 & 965 \\
Active hard words & 43\,036 & 43\,028 \\
\midrule
Secondary decoder epoch & \multicolumn{2}{c}{140} \\
Decoder MSE & 0.01134 & 0.01121 \\
Decoder PSNR (dB) & 19.45 & 19.50 \\
Decoder MS-SSIM & 0.714 & 0.718 \\
\bottomrule
\end{tabular}
\caption{DINO-defined relational learning on Flowers102. Representation and secondary decoder checkpoints use separate validation criteria. The encoder is frozen during decoder training. The study uses one seed.}
\label{tab:flowers102-teacher-results}
\end{table}

On test images, teacher-favored pairs collide with probability $1.29\times10^{-2}$, compared with aggregate collision $1.04\times10^{-3}$.
Teacher-disfavored pairs collide with probability $1.47\times10^{-4}$.
The corresponding enrichment is $12.4\times$ and suppression is 85.8\%.
The weighted mean teacher cosine is 0.579 on favored relations and 0.311 on disfavored relations; both are positive, and their separation confirms that the baseline-relative split behaves as intended.

The hard representation has about 965 collision-effective words across 43\,028 active words, reflecting broad support with concentrated occupancy.
The result shows that pair-specific alignment and relational separation can organize a noncollapsed code without an explicit marginal entropy penalty.
The evaluation covers the declared relational distinctions and their hard-collision diagnostics rather than numerical reconstruction of the full teacher geometry.

\begin{figure}[t]
\centering
\includegraphics[width=0.68\linewidth]{flowers102_teacher_graph_dynamics.png}
\caption{Validation dynamics of the baseline-relative teacher study. Hard collisions on favored pairs are enriched and those on disfavored pairs are suppressed relative to the pooled occupancy baseline as the code develops. The secondary decoder curve belongs to a separate stage after the encoder is frozen.}
\label{fig:flowers102-teacher-dynamics}
\end{figure}

After representation learning, a decoder with the synthesis-transform form of Study 4 is trained on frozen hard codes and selected by validation MSE.
It achieves test PSNR 19.50 dB and MS-SSIM 0.718.
This shows that the retained representation supports nontrivial image reconstruction even though pixel fidelity did not shape it directly.
In the absence of random-encoder and shuffled-teacher decoder controls, the result serves as a compatibility probe between two fidelities rather than a causal attribution or comparison with the reconstruction-trained codec.
Qualitative examples are retained in \cref{fig:flowers102-teacher-reconstruction-examples}.

The learned code remains noncollapsed while placing substantially more collisions on teacher-favored pairs and fewer on teacher-disfavored pairs than the pooled occupancy baseline, demonstrating that pair-specific relational requirements can directly organize finite-codeword equality without a decoder in representation learning.

